% docs/manuscript/sections/protocol_v2.tex
% Experimental Protocol -- stream S7-ter, LOT A. Answers reviewer #3's nominative demand for the
% resampling unit of every confidence interval, the inter-seed aggregation rule, the treatment of
% the 12 ProteuS transitions, the complete hyper-parameter set, the multiplicity policy, the
% complete separations, the legitimacy of the seed-paired bootstrap and the per-worker PRNG lock.
% Every value below is quoted from a committed artifact or from a source line, never from a report.

\section{Experimental Protocol}\label{sec:protocol}

This section states, for every interval and every test reported in
Section~\ref{sec:experiments}, what was resampled, what was pooled and what was held fixed. The
answer is not uniform across the paper, and it is not presented as if it were.

\subsection{Resampling Unit of Each Confidence Interval}\label{sec:protocol_ci}

Five distinct resampling units appear in the paper. Table~\ref{tab:protocol_ci} names each one
with its replicate count, its seed and its generator.

\begin{table*}[t]
  \centering
  \caption{Resampling unit of every confidence interval reported in this paper. The unit is not
  uniform: it is the seed wherever arms are seed-paired, the per-drift jump where the estimand is
  itself a per-drift quantity, and parametric where the null is a fitted family. Every generator is
  a locally injected \texttt{numpy.random.Generator}; no bootstrap draws from the global NumPy
  state.}
  \label{tab:protocol_ci}
  \begin{tabular}{@{}llrll@{}}
    \toprule
    \textbf{Statistic} & \textbf{Resampling unit} & \textbf{$n_{\mathrm{boot}}$} & \textbf{Seed} & \textbf{Generator} \\
    \midrule
    Table~\ref{tab:results_concept} F1 and ADD; $\alpha$-sweep & seed index, paired across arms & 1000 & 12345 & local \texttt{default\_rng} \\
    Hydra RMST ratio $\tau_{\mathrm{HAT}}/\tau_{\mathrm{ARF}}$ & seed index, paired across arms & 10000 & 20260906 & local \texttt{default\_rng} \\
    Evidence ceiling, $\lambda_{\mathrm{op}}$, $\Delta e_c$ & seed, paired and unpaired variants & 10000 & 12345 & local \texttt{default\_rng} \\
    $\Delta e$ on real streams (Table~\ref{tab:real_data_summary}) & the per-drift jump, $K$ per variant & 1000 & 42 & local \texttt{default\_rng} \\
    Exponential-fit KS $p$-value for $\tau_{\mathrm{HAT}}$ & parametric, under the fitted exponential & 2000 & $\lfloor 1000\,\Delta e \rfloor$ & local \texttt{default\_rng} \\
    \bottomrule
  \end{tabular}
\end{table*}

\begin{remark}[Correction to the submitted resampling unit of Table~\ref{tab:results_concept}]\label{rem:boot_unit}
  The submitted version resampled the 360 \emph{runs} of a cell with replacement, off the global
  NumPy state seeded once before aggregation, while the caption asserted the seed as the unit of
  statistical independence. The two statements are incompatible: the 36 streams a seed produces
  share that seed's forest initialisation and are not independent of one another, so a run-level
  resample understates the width of the interval. The bootstrap now resamples the 30 seed indices,
  paired across arms, on a locally injected generator, which is the scheme the Hydra ratio already
  used. The caption's independence claim and the computation now agree.
\end{remark}

Two degeneracies are reported as such rather than as tight intervals. The $\Delta e$ estimate for
\emph{gradual\_balanced} rests on a single canonical drift ($K = 1$), so every bootstrap replicate
returns the same value and the interval collapses to the point estimate $0.450341$; it carries no
information about sampling variability. The detection-rate contrast between KSWIN$+$ARF and
PHT$+$ARF at $c_{\mathrm{int}} = 1$ separates at every one of the 30 seeds with no seed-level
dispersion at all ($1080/1080$ against $0/1080$, per-seed rates $1.00$ against $0.00$), so its
seed-level interval is likewise of zero width.

\paragraph{A non-adapting control for the error-jump estimate.}
The effective error jump on the real streams is read from the error trace of a reference tree that
learns at every step, which is the same adaptation the quantity is supposed to expose: the estimator
is measured on a model already absorbing the change. Every reported value is therefore paired with
an oracle computed in the same pass, from a copy of that tree forked once at the end of the warm-up
and never trained again --- predictions only, no updates. Oracle and estimate share the drift
positions, the adaptive windows and the bootstrap seed, so they differ in exactly one input.

Reading the oracle requires reading the frozen model's own error level beside it, and the artifact
carries it for that reason. A frozen model at chance cannot exhibit an error jump, so an oracle near
zero means ``no jump in the stream'' only while the frozen model still discriminates; otherwise it
means ``nothing left to lose'', and the two are indistinguishable from the jump alone.

\subsection{Aggregation Across Seeds}\label{sec:protocol_aggregation}

A run is one $(\text{transition}, \text{GARCH regime}, \text{seed})$ triple evaluated by one
pipeline. A run counts as a \emph{detection} when its $\mathrm{F1}$ is strictly positive, which on
the tolerance-window protocol of Section~\ref{sec:proteus} is equivalent to at least one alarm
falling inside $[\tau^*, \tau^* + \tau]$. Detections are summed \emph{within} a seed over the 36
streams that seed produces ($12$ transitions $\times$ $3$ GARCH regimes), giving one count per seed
per pipeline. Two pipelines are compared by pairing those counts on the seed and applying a
two-sided sign test to the 30 paired differences; ties contribute to neither side and reduce
$n_{\mathrm{eff}}$. Seeds are never pooled across transitions, and no test is ever computed on the
1080 runs treated as independent units, which is the pseudo-replication the comparison literature
warns against~\cite{demsar_statistical_2006} --- the run-level Fisher figure recorded alongside each pair
in the artifact is retained for documentation only and is not reported as evidence.

Cell means in Table~\ref{tab:results_concept} are plain means over the 360 runs of a
$(\text{detector}, \text{clock}, \text{regime})$ cell; only the interval around them is computed at
the seed level. $\mathrm{ADD}$ is averaged over detected runs alone and is undefined --- printed
$\infty$ --- when a cell detects nothing, which is the reading of the three blind-spot rows.

\subsection{The Twelve ProteuS Transitions}\label{sec:protocol_transitions}

The 12 transitions are the ordered pairs of the four reference ETFs, and the drift width is a
function of the pair rather than a crossed factor: $w = 100$ for the three transitions \emph{into}
VNQ, $w = 1000$ for the three \emph{out of} it, and $w = 500$ for the remaining six. The evaluation
grid is the Cartesian product of the 12 transitions with the 30 seeds; the three GARCH regimes are
evaluated inside each grid point on streams regenerated from the same seed. Hence $360$ runs per
table cell ($12 \times 30$) and $1080$ per pipeline across the three regimes
($30 \times 36$). Each transition keeps its own tolerance $\tau = \max(1000, w)$, so the scoring
window follows the drift width and is never averaged across widths.

Transitions are a within-seed factor: they are pooled inside a seed to form that seed's detection
count, and the seed remains the unit carried into every test. Pooling in the other order --- across
seeds within a transition --- would treat 30 draws of the same stream family as replicates of the
transition and is not done anywhere in this paper.

\subsection{Detector and Ensemble Hyper-Parameters}\label{sec:protocol_hyper}

Table~\ref{tab:protocol_hyper} gives the complete parameterisation. Every value is pinned in the
single source of truth \texttt{config/experiment\_ssot.py} or in the factory functions of the
experiment under test; none is a local literal, and a constant re-bound to a literal fails the
repository's consistency test without importing any experiment.

\begin{table*}[t]
  \centering
  \caption{Complete detector and ensemble parameterisation. River 0.23.0 throughout. Parameters not
  listed are River defaults. $c$ denotes the internal ADWIN clock of the adaptive forest, swept over
  $\{1, 32\}$; $c_{\mathrm{ext}}$ the clock of an external ADWIN monitor.}
  \label{tab:protocol_hyper}
  \begin{tabular}{@{}lll@{}}
    \toprule
    \textbf{Component} & \textbf{Parameterisation} & \textbf{Where used} \\
    \midrule
    Adaptive Random Forest & $M = 10$; \texttt{drift\_detector} $=$ ADWIN($\delta = 0.002$, clock $= c$); & all adaptive arms \\
                           & \texttt{warning\_detector} $=$ ADWIN($\delta = 0.002$, clock $= c$) & \\
    Hoeffding Tree (non-adaptive baseline) & River defaults, no reset & baseline arms \\
    Static Random Forest & bagging of $M = 10$ Hoeffding Trees, no internal detector & Section~\ref{sec:solution_rf} \\
    Streaming Random Patches & $M = 10$, subspace $0.6$, patches, $\lambda = 6.0$; both detectors ADWIN($0.002$, $c$) & Section~\ref{sec:proteus} \\
    \midrule
    Page--Hinkley (external) & $\lambda = 15$, $\delta_P = 0.005$ (ProteuS); $\lambda = 25$, $\delta_P = 0.005$ (crossover) & Sections~\ref{sec:proteus},~\ref{sec:crossover} \\
    Strict CUSUM (external) & $\delta_P = 0.01$, $\lambda \in \{2.5, 5, 8, 10, 15, 20, 25, 50, 100\}$; & Section~\ref{sec:bernoulli_validation} \\
                            & $p_0$ from the 1000-step pre-drift warm-up (see below) & \\
    ADWIN (external) & $\delta = 0.002$, clock $= c_{\mathrm{ext}}$ & Sections~\ref{sec:proteus},~\ref{sec:hardware} \\
    KSWIN (external) & $\alpha \in \{0.001, 0.005, 0.01, 0.05\}$, window $100$, statistic $30$, & Section~\ref{sec:proteus} \\
                     & applied to the error mean over a 30-step buffer & \\
    EDDM (external) & River defaults & Section~\ref{sec:proteus} \\
    \midrule
    ProteuS streams & $n = 8000$, $\tau^* = 4000$, $w \in \{100, 500, 1000\}$, Student-$t_7$ innovations & Section~\ref{sec:proteus} \\
    Bernoulli crossover streams & $n = 8000$, $\tau^* = 4000$, tolerance $1000$, $\Delta e$ on a 15-point grid & Section~\ref{sec:crossover} \\
    Real streams & warm-up $\min(10^5, 0.10\,n)$, tolerance $\min(5000, 0.05\,n)$ & Table~\ref{tab:real_data_summary} \\
    \bottomrule
  \end{tabular}
\end{table*}

\paragraph{The reference level $p_0$ of the external CUSUM.}
The cumulative monitor of Section~\ref{sec:bernoulli_validation} tests the error stream against a
pre-change level $p_0$, and how that level is obtained is part of the protocol rather than an
implementation detail: it sets the drift term of the accumulation, so a $p_0$ above the true
pre-change error rate adds slack and starves the monitor earlier, while one below it accumulates on
noise. Two constructions coexist in the race-condition experiment, on the same error stream, at the
same tolerance and the same threshold, differing only in $p_0$:

\begin{itemize}
  \item a \emph{nominal} arm with $p_0 = 0.05$ fixed before the stream starts --- the historical
        value, retained so the original configuration stays reproducible;
  \item an \emph{empirical} arm with $p_0$ set at the change point to the mean error over the
        preceding 1000 steps.
\end{itemize}

\textbf{Every metric this paper reports for that experiment is the empirical arm.} The blind-spot
indicator and the detection rate are both defined on the empirical stopping time, and Figure~1 plots
it; the nominal arm is a secondary control column, recorded in the artifact and published nowhere.

The distinction is not cosmetic, and it is worth decomposing rather than asserting. At
$\lambda = 25$ the reported blind-spot share of $0.895$ is $0.815$ of runs in which adaptation
pre-empts the monitor --- both stopping times finite, the internal swap first --- plus $0.080$ in
which the monitor never fires at all within the horizon; the reported detection rate is the
complement of that second term, $0.920$. Under the nominal arm the same decomposition reads $0.355$
pre-empted plus $0.610$ censored for a share of $0.965$, at a detection rate of $0.390$. The two
arms therefore agree that the pipeline is blind and disagree completely about \emph{why}: with
$p_0 = 0.05$ fixed well above the measured pre-drift error rate, the nominal monitor carries so much
slack that it is starved outright in six runs out of ten, and the pre-emption the experiment is
built to measure all but disappears behind that. Quoting a figure from one arm as the other's would
misstate not only its value but its mechanism.

Every later campaign --- the instrumented starvation experiments, the real-world streams, the
$\lambda_{\mathrm{op}}$ sweep and the synchronised traces --- calibrates $p_0$ on the warm-up and
carries one arm only. The fixed literal survives in those as an unreachable fallback for an empty
warm-up buffer, which a 1000-step window never produces.

\paragraph{The warning detector is two parameters, not one.}
River~0.23.0~\cite{montiel_river_2021} resolves an unset \texttt{warning\_detector} on its adaptive forest to
ADWIN($\delta = 0.01$, clock $= 32$) and an unset \texttt{drift\_detector} to
ADWIN($\delta = 0.001$, clock $= 32$); a bare \texttt{ADWIN()} is neither, at
$\delta = 0.002$, clock $= 32$. The submitted version pinned only the drift detector on the ProteuS
and crossover arms, so their warning detector ran at $\delta = 0.01$ and clock $32$ while the
instrumented arms pinned both at $\delta = 0.002$ and clock $c$. Unifying the two therefore moves
\emph{two} parameters, $0.01 \to 0.002$ and $32 \to c$, not the clock alone.

Both consequences of that unification are stated, because they coexist and neither cancels the
other. The warning configuration is mechanically \emph{determining}: it gates when a background
tree starts training, hence how much a replacement has learned when it is installed, hence the
duration of the post-drift transient the external monitor sees. And wherever the two detectors are
pinned identically it is simultaneously \emph{inert}: they become identically parameterised clones
fed the same error sequence, so they fire at the same step, the background tree is created and
promoted within the same update, and nothing ever trains in between. Instrumentation on the pinned
build measures this directly, at both clocks: warning and drift fire element-wise together, the
maximum number of simultaneously live background trees is zero, and \emph{every} replacement
installs a tree created in the same step and warmed by nothing --- 47 of 47 at $c = 1$ and 36 of 36
at $c = 32$, with 0 warmed background trees in both. A mechanism that determines the transient and
never runs is not a contradiction; it is why the effect is insensitive to the warning threshold in
the tested configuration, and why it would not be under a warning detector that genuinely leads the
drift detector --- which is what River's own defaults provide, and what the control arm of
Section~\ref{sec:protocol_arms} measures.

\subsection{Control Arm: River's Own Defaults}\label{sec:protocol_arms}

Unifying the warning detector onto the drift detector buys comparability across experiments at the
cost of disabling the one mechanism that would otherwise give a replacement tree something to learn
before it is installed. Publishing the phenomenon only under that configuration would leave it
measured on a forest whose background-tree mechanism is structurally inert, and a reader who knows
the library would be right to ask whether the blind spot is a property of adaptation or of a
crippled instance of it.

Both arms are therefore measured and both are reported:

\begin{itemize}
  \item \textbf{Unified} --- \texttt{warning\_detector} $=$ \texttt{drift\_detector} $=$
        ADWIN($\delta = 0.002$, clock $= c$). Comparable across every experiment in the paper; the
        warning fires with the drift and no background tree ever trains.
  \item \textbf{Library defaults} --- \texttt{warning\_detector} $=$ ADWIN($\delta = 0.01$,
        clock $= 32$), the drift detector unchanged. This is what a practitioner who instantiates
        the forest without naming a warning detector actually deploys, and the only configuration in
        which the warning genuinely leads the drift.
\end{itemize}

Both arms were run in full and are reported below. The measurement did not go the way the question
was posed: the phenomenon is \emph{stronger} under the library defaults, not weaker.

\paragraph{Heteroscedastic streams: invariant where the claims live.}
On the ProteuS evaluation the unification changes nothing the paper asserts. The three blind-spot
collapses stay at $\mathrm{F1} = 0.00$ exactly with $\mathrm{ADD}$ undefined in all three GARCH
regimes; the KSWIN resolution stays at $\mathrm{F1} = 1.00$, $\mathrm{ADD} = 14\,[0]$ with zero
variance at both clocks; the $\alpha$-sweep is invariant down to the raw per-run records. Rows
whose classifier carries no warning detector --- the Hoeffding Tree, the static forest, and the
patches ensemble, which already pinned both --- are bit-identical. What moves is confined to the
safe-zone rows of the adaptive forest at $c = 32$ (F1 $0.63 \to 0.67$ and $0.56 \to 0.61$) and to a
single seed-level count, the windowed monitor's $1063/1080 \to 1070/1080$.

\paragraph{Controlled Bernoulli streams: not invariant, and the direction matters.}
On the regime crossover the two arms separate, and only on the adaptive arm --- the fifteen rows of
each non-adaptive pipeline are identical under both, as they must be.

\begin{table}[t]
  \centering
  \caption{Regime crossover, adaptive arm: library defaults against the unified configuration.
  Missed-detection rate in percent over 100 streams per point. The non-adaptive pipelines are
  identical under both arms and are omitted.}
  \label{tab:protocol_arm_r3}
  \begin{tabular}{@{}lrr@{}}
    \toprule
    \textbf{Quantity} & \textbf{Library defaults} & \textbf{Unified} \\
    \midrule
    Missed detections at $\Delta e = 0.50$ & 100 & 80 \\
    Missed detections at $\Delta e = 0.47$ & 98  & 65 \\
    Missed detections at $\Delta e = 0.26$ & 1   & 4  \\
    Missed detections at $\Delta e = 0.16$ & 1   & 13 \\
    Missed detections at $\Delta e = 0.09$ & 1   & 69 \\
    \midrule
    Post-drift accuracy at $\Delta e = 0.50$ & 0.984 & 0.980 \\
    Accuracy gap against the static forest (pp) & 23.41 & 22.99 \\
    Pre-drift false-alarm ratio, non-adaptive over adaptive & 1.98 & 2.00 \\
    \bottomrule
  \end{tabular}
\end{table}

Two readings follow, and both are consequences of the same mechanism. At the top of the band the
library defaults starve the monitor completely while the unified configuration leaves one run in
five detectable: a warning detector that genuinely leads the drift detector installs a background
tree that has trained, the post-drift transient is shorter, and the cumulative monitor accumulates
less. In the safe zone the ordering reverses --- at $\Delta e = 0.09$ the unified arm misses $69\%$
against $1\%$ --- because an inert warning mechanism also removes the noise-triggered replacements
that were producing early alarms there.

The consequence is stated rather than absorbed. The claim that missed detections reach $100\%$ at
$\Delta e = 0.50$ holds on the library-defaults arm and \emph{does not survive} the unification,
where the figure is $80\%$. What survives on both arms is the qualitative claim --- the adaptive
pipeline starves across the band while both non-adaptive pipelines detect every drift --- together
with the accuracy gap, which reads $23.4$ against $23.0$ percentage points, and the false-alarm
halving, which reads $1.98$ against $2.00$.

The arm a figure comes from is therefore part of its statement, and this section names it rather
than leaving it to be inferred. \textbf{Every regime-crossover figure in this paper is the unified
arm}, chosen so that a single warning configuration holds across every experiment reported here; the
library-defaults arm is retained as the comparability control and is reproducible byte for byte from
the artifact repository.

The cost of that choice is stated, not absorbed: the phenomenon is published under the weaker of the
two configurations measured. Under the library defaults the cumulative monitor is starved outright
at the top of the band, and the figure the submitted version quoted --- $100\%$ --- is that arm's.
Reporting the unified arm is therefore conservative with respect to the paper's own claim, which is
the only direction in which a methodological choice may err.

\subsection{Multiplicity}\label{sec:protocol_multiplicity}

The family of hypothesis tests reported in this paper is stated in full: six seed-level pipeline
contrasts on ProteuS, one contrast in the KSWIN $\alpha$-sweep, and three seed-level contrasts on
the INSECTS variants of Table~\ref{tab:real_data_summary} --- ten tests. The three BAF rows carry
no test, their $\Delta e$ being indistinguishable from zero (Section~\ref{sec:protocol_ci}).

Eight of the ten $p$-values sit exactly at the resolution floor of the two-sided sign test at
$n_{\mathrm{eff}} = 30$, derived in Section~\ref{sec:protocol_separation}. No correction over a
family of this size can lift them: Holm--Bonferroni at $\alpha = 0.05$ compares the smallest
against $\alpha/10 = 5 \times 10^{-3}$, and $2^{-29} \approx 1.86 \times 10^{-9}$ clears that by six
orders of magnitude, as it clears every subsequent step. Bonferroni over the whole family gives
$10 \cdot 2^{-29} \approx 1.86 \times 10^{-8}$, still far below any conventional level. For these
eight the correction is not applied because it is arithmetically incapable of changing a decision,
which is a justification and not an omission.

The two remaining tests are reported with the opposite verdict, and it is the honest one. The
KSWIN $\alpha$-sweep contrast is degenerate ($p = 1$: both arms detect 1080 of 1080 runs), and the
INSECTS \emph{abrupt\_balanced} contrast has $p = 0.043$, which under Holm is compared against
$\alpha/2 = 0.025$ and \textbf{does not survive}. We therefore do not claim a seed-level effect on
\emph{abrupt\_balanced}. That is consistent with how Section~\ref{sec:proteus} already treats that
variant --- a weak witness whose five canonical jumps are heterogeneous and partly negative, and
whose marginal advantage is explicitly not attributed to the mechanism. The multiplicity correction
removes a claim the paper had already declined to make.

\subsection{Complete Separations and Effect Sizes}\label{sec:protocol_separation}

Several contrasts separate completely: one arm detects in every seed, the other in none. The
$p$-value reported for these is not an estimate and must not be read as one. With all 30 seeds
favouring the same arm, the two-sided sign test returns
$2 \cdot 2^{-30} = 2^{-29} = 1.8626451 \times 10^{-9}$ exactly --- the smallest value the test can
produce at $n_{\mathrm{eff}} = 30$, attained whenever separation is complete and independent of how
large the separation is. We therefore report it as $p \le 2^{-29}$ and pair every such contrast
with an effect size, since the $p$-value carries no magnitude information whatsoever.

The effect size is the difference in detection rate between the two arms, computed per seed over
that seed's 36 streams and then averaged over the 30 seeds, so that it is commensurable with the
test that accompanies it. Table~\ref{tab:protocol_effects} reports it with a 95\% interval taken
across seeds.

\begin{table*}[t]
  \centering
  \caption{Effect sizes for the seed-level contrasts of Section~\ref{sec:experiments}. $\Delta$ is
  the mean over 30 seeds of the within-seed detection-rate difference (36 streams per seed), in
  percentage points, with its 95\% seed-level interval. Every contrast listed separates completely
  and reports $p \le 2^{-29}$; the $p$-value is identical across rows by construction, the effect
  size is not.}
  \label{tab:protocol_effects}
  \begin{tabular}{@{}llrrr@{}}
    \toprule
    \textbf{Arm A} & \textbf{Arm B} & \textbf{A} & \textbf{B} & \textbf{$\Delta$ (pp, 95\% CI)} \\
    \midrule
    PHT $+$ ARF ($c{=}1$)   & PHT $+$ HT              & 0/1080    & 932/1080  & $-86.30\ [-87.48, -85.11]$ \\
    EDDM $+$ ARF ($c{=}1$)  & EDDM $+$ HT             & 0/1080    & 882/1080  & $-81.67\ [-82.90, -80.43]$ \\
    SRP $+$ PHT ($c{=}1$)   & PHT $+$ HT              & 0/1080    & 932/1080  & $-86.30\ [-87.48, -85.11]$ \\
    KSWIN $+$ ARF ($c{=}1$) & PHT $+$ ARF ($c{=}1$)   & 1080/1080 & 0/1080    & $+100.00$ (degenerate) \\
    PHT $+$ RF (static)     & PHT $+$ ARF ($c{=}1$)   & 913/1080  & 0/1080    & $+84.54\ [83.33, 85.75]$ \\
    ADWIN $+$ RF (static)   & ADWIN $+$ ARF ($c{=}1$) & 959/1080  & 1070/1080 & $-10.28\ [-11.87, -8.69]$ \\
    \bottomrule
  \end{tabular}
\end{table*}

The last row is the one that would be misread from its $p$-value alone. It carries the same
$p \le 2^{-29}$ as the blind-spot collapses, because all 30 seeds favour the same arm, yet its
effect is an order of magnitude smaller: the static forest detects $10.3$ percentage points
\emph{less} than the adaptive one under a windowed monitor. Reported as a $p$-value it looks like
the collapses; reported as an effect size it is what it is, a small and consistent disadvantage,
and it is the measurement behind the statement that the static ensemble resolves the starvation
regime rather than replacing adaptation in general.

\subsection{Legitimacy of the Seed-Paired Bootstrap}\label{sec:protocol_paired}

Pairing on the seed is legitimate only if the arms see the same stream at the same seed. On the
instrumented Bernoulli family this holds by construction and is verifiable in one line of each
experiment: the covariate stream is produced by exactly two scalar draws per step,
\texttt{x0, x1 = rng.normal(), rng.normal()}, from a generator seeded by the run's own seed, inside
the step loop and before any model call. Since neither \texttt{predict\_one} nor \texttt{learn\_one}
consumes from that generator, the sequence $(x_{0,t}, x_{1,t})_{t<n}$ is bit-identical across arms
for a given (boundary shift, seed) pair, and the label rule $y_t = \mathbf{1}[x_{0,t} + x_{1,t} > b]$
is a deterministic function of it. Arms at equal seed therefore differ only in the classifier and
the monitor, which is exactly the contrast the paired resampling estimates.

Two deviations from that invariant exist in the repository and are declared rather than smoothed
over. The regime-crossover experiment draws its covariates from the legacy
\texttt{RandomState} (MT19937) generator rather than the modern one, deliberately, to reproduce the
streams of the submitted version bit-for-bit; pairing is unaffected, since both arms read the same
legacy stream, but its draws are not comparable with those of the other experiments. And the
$\lambda_{\mathrm{op}}$ sweep seeds only its local generator, without the global lock described
next; it reports no paired interval, so nothing in the paper depends on that omission.

\subsection{Per-Worker Pseudo-Random Number Generation}\label{sec:protocol_prng}

Every parallel worker of the instrumented experiments opens with the same three statements, in this
order:

\begin{center}
  \begin{tabular}{@{}l@{}}
    \texttt{safe\_seed = int(seed \% (2**31 - 1))} \\
    \texttt{random.seed(safe\_seed)} \\
    \texttt{np.random.seed(safe\_seed)} \\
    \texttt{rng = np.random.default\_rng(safe\_seed)} \\
  \end{tabular}
\end{center}

The local generator \texttt{rng} produces the covariate stream. The two global seedings are not
redundant and must not be removed as cleanup: River's Cython tree-spawn path draws the Poisson
bagging weights and the split randomisation from the global NumPy singleton and from the standard
library's global generator, neither of which accepts an injected generator through the public API.
Without the global lock the forest's internal stochasticity varies with worker scheduling and the
artifacts cease to be bit-reproducible. The modulo is not cosmetic either: the Cython extensions
accept a 32-bit signed seed, and an unreduced 64-bit value overflows silently.

The ProteuS experiments carry a variant of the same lock --- \texttt{random.seed} and
\texttt{np.random.seed} on the reduced seed, with no local generator at all --- because their
covariate stream is generated by NumPy's global state inside the stream simulator rather than by an
injected generator. The lock is therefore load-bearing there for the streams themselves, not only
for River's internals.

All experiments run under \texttt{PYTHONHASHSEED=0}, set by the execution wrappers, so that
dictionary iteration order inside River's feature maps is fixed. The host specification, the pinned
package versions and the measured per-stage wall-clock times are recorded in the artifact
repository.
